In this chapter we will present the classical Calogero-Moser systems under consideration and some of their properties. We will show how they can be arrived at by means of the Kazhdan-Kostant-Sternberg (KKS) construction and how solutions can be obtained through the Olshanetsky-Perelomov (OP) projection method. Lastly, we will report their Lax representations and their Bäcklund transformations, from which we will derive soliton equations in Chapter~\ref{ch:solitons}.


\section{Calogero-Moser (CM) systems}

\begin{defi}
Let $(q, p) = (q_1, \dots, q_N, p_1, \dots, p_N)$ be a set of canonical coordinates. A dynamical system is said to be an $A_{N-1}$ Calogero-Moser \textnormal{(CM)} system with some external potential $W(x)$ if it is characterized by the Hamiltonian

\begin{equation}
\mathcal{H}_{CM} = \sum_{i=1}^N \Big(\frac{1}{2} p_i^2 + W(q_i)\Big) + g^2 \sum_{1 \leq i < j \leq N} V(q_i - q_j)\;,
\end{equation}

(with unit masses), where $g^2$ is the coupling constant and the interparticle potential $V (x)$ belongs to the classified set of meromorphic functions

\[V(x) = \left \{ \begin{array}{ccc}
{1}/{x^2} & \textnormal{(rational)} & \textnormal{(I)}\\
{a^2}/{\sin^2{ax}} & \textnormal{(trigonometric)} & \textnormal{(II)}\\
{a^2}/{\sinh^2{ax}} & \textnormal{(hyperbolic)} & \textnormal{(III)}\\
\wp (x) & \textnormal{(elliptic)} & \textnormal{(IV)}
\end{array}
\right.\]

where $\wp (x)$ is the \textit{elliptic, doubly periodic Weierstrass function}

\begin{equation}
    \wp (z, \omega_1, \omega_2) = \frac{1}{z^2} + \sum_{m, n \neq (0,0)} \Big (  \frac{1}{(z - 2m\omega_1 - 2n\omega_2)^2} - \frac{1}{(2m\omega_1 + 2n\omega_2)^2}\Big )
\end{equation}

with periods $2\omega_1$ and $2\omega_2$.
\end{defi}

These forms of the potential are not arbitrary. They are fixed by the form of the Lax pair proposed by F. Calogero (see Section~\ref{laxcouples}). The equivalence between the equations of motion and the Lax evolution equation puts strong constraints on Calogero's ansatz, which result in the functional equation (\ref{laxfunctionalequations}) for the Lax function. Its general solution is given in terms of the Weierstrass elliptic function $\wp(x)$ (see \cite{calogero2003classical} for a detailed derivation). In fact, the type I-III systems can be viewed as limiting cases of the type IV system \cite{ruijsenaars1999systems}:

\begin{equation}
    \wp(x) = \left \{ \begin{array}{ccc}
    {1}/{x^2} & \omega_1 = \infty & \omega_2 = \infty\\
    {a^2}/{\sin^2{ax}} - {a^2}/{3} & \omega_1 = {\pi} / {2a} & \omega_2 = i\infty\\
    {a^2}/{\sinh^2{ax}} + {a^2}/{3} & \omega_1 = \infty\ & \omega_2 = {i\pi}/{2a}
\end{array}
\right.
\end{equation}

More generally, the potential $V(x)$ is constructed from root systems of Lie algebras, as reported in \cite{perelomov1989}. \\
From Hamilton's equations we get the equation of motion of the $i^{\textnormal{th}}$ particle

\begin{equation}
    \ddot q_i = - W'(q_i) - g^2 \sum_{j = 1,\;j \neq i}^N V'(q_i - q_j)\;,
\end{equation}

where we used that $V$ is even. In this work, we will deal with the classical CM systems I, II and III. We will also consider the rational Calogero-Moser system with external harmonic potential, $W(x) = \frac{1}{2}\omega^2 x^2$, denoted by $\textnormal{I}^+$ or eCM.

\section{KKS construction of CM systems}
Given a dynamical system with a Hamiltonian $\mathcal{H}:\mathcal{T}^\ast \mathcal{M} \to \mathbb{R}$ invariant under a continuous 1-parametric symmetry group, Noether’s theorem \cite{noether1918} gives rise to integrals of motion corresponding to this symmetry. The procedure of \textit{Hamiltonian reduction} consists in eliminating some degrees of freedom by setting these integrals to constant values. In this way we are able to obtain a new submanifold, $\mathcal{T}^\ast \mathcal{M}_{red}$ of lower dimension than the initial one, usually simplifying the problem\footnote{For a more formal treatment, see V.~I. Arnold, V.~V. Kozlov and A.~I. Neishtadt, \textit{Mathematical Aspects of Classical and Celestial Mechanics}, 3rd ed., Encyclopaedia of Mathematical Sciences, Springer, Berlin, 2007, in particular Section 3.2.2, pp.~116--120.}.
In this case, the reduced Hamiltonian typically is more complicated than the original one, suggesting in reverse to construct a system with more degrees of freedom where the hidden symmetries are apparent. \\
In the Kazhdan-Kostant-Sternberg (KKS) construction of CM systems \cite{KKS1978}, CM systems emerge from Hamiltonians on the larger configuration space $X^0_N = \{X\}$ of $N \times N$ Hermitian matrices $X$. We set the phase manifold as the cotangent bundle $\mathcal{T}^\ast X^0_N$, mapped by canonical pairs of $N \times N$ matrices $(X,Y)$. The new system is chosen to be dynamically simple, so that its equation of motion can be solved in closed form \cite{perelomov1989}. \\
If $U$ is a unitary $N \times N$ matrix, then the Hamiltonians are invariant with respect to similarity transformations

\begin{equation}
    (X, Y) \mapsto (UXU^{-1}, UYU^{-1})\;.
\end{equation}

Let us introduce the following $N \times N$ Hermitian matrices:

\begin{equation}
\begin{cases}
Q_{nm} = \delta_{nm} q_n(t) \\
L_{nm} = \delta_{nm}p_n(t) + (1 - \delta_{nm}) \frac{ig}{q_n - q_m}\;.
\end{cases}
\end{equation}

Such matrices are made to correspond to generic Hermitian matrices $X, Y \in X_N^0$ (see \cite{olshanetskyperelomov} for details) by a similarity transformation that leaves the Hamiltonian of the system invariant:

\begin{equation}
    X = UQU^{-1}\;\;\;\;,\;\;\;\;Y = ULU^{-1}\;.
\end{equation}

We report below the construction for each of the cases under analysis.

\subsection{Rational CM (rCM) system}
The KKS construction of the rational CM system leads to the Hamiltonian

\begin{equation}
    \mathcal{H}_{rCM} = \frac{1}{2}\operatorname{tr}(Y^2), \;\;\;\;\;\;(X,Y) \in \mathcal{T}^\ast X^0_N
\end{equation}

describing free motion on $X^0_N$. Applying the similarity transformation we get

\begin{equation}
    \mathcal{H}_{rCM} = \frac{1}{2}\operatorname{tr}(Y^2) = \frac{1}{2}\operatorname{tr}(UU^{-1}Y^2) = \frac{1}{2}\operatorname{tr}(U^{-1} Y^2 U) = \frac{1}{2}\operatorname{tr}(L^2)\;.
\end{equation}

It is easy to verify that the projection $X \to Q$, $Y \to L$ reduces the system to the rCM one:
    
\begin{equation}
    \mathcal{H}_{rCM} = \frac{1}{2}\operatorname{tr}(L^2) = \frac{1}{2}\sum_{i=1}^N p_i^2 + g^2 \sum_{1 \leq i < j \leq N} \frac{1}{(q_i - q_j)^2}\;.
\end{equation}

\subsection{Rational CM system with external harmonic potential (eCM)}
The KKS construction of the rational CM system with external harmonic potential starts from the Hamiltonian

\begin{equation}
    \mathcal{H}_{eCM} = \frac{1}{2}\operatorname{tr}(Y^2 + \omega^2 X^2), \;\;\;\;\;\;(X,Y) \in \mathcal{T}^\ast X^0_N\;,
\end{equation}

describing harmonic motion on $X^0_N$.

Applying the same similarity transformation we get

\begin{equation}
\begin{split}
    \mathcal{H}_{eCM} &= \frac{1}{2}\operatorname{tr}(Y^2) + \frac{1}{2}\omega^2 \operatorname{tr}(X^2) = \\  &= \frac{1}{2}\operatorname{tr}(UU^{-1}Y^2) + \frac{1}{2}\omega^2 \operatorname{tr}(UU^{-1} X^2) = \\ 
    & = \frac{1}{2}\operatorname{tr}(U^{-1} Y^2 U) + \frac{1}{2}\omega^2 \operatorname{tr}(U^{-1} X^2 U)= \frac{1}{2}\operatorname{tr}(L^2) + \frac{1}{2}\omega^2 \operatorname{tr}(Q^2)\;.
\end{split}
\end{equation}

The projection then gives
    
\begin{equation}
    \mathcal{H}_{eCM} = \frac{1}{2}\operatorname{tr}(L^2) + \frac{1}{2}\omega^2 \operatorname{tr}(Q^2)= \frac{1}{2}\sum_{i=1}^N (p_i^2 + \omega^2 q_i^2) + g^2 \sum_{1 \leq i < j \leq N} \frac{1}{(q_i - q_j)^2}\;.
\end{equation}


\subsection{Trigonometric CM (tCM) system}
The trigonometric CM system derives from the Hamiltonian on $\mathcal{T}^\ast X^0_N$

\begin{equation}
    \mathcal{H} = \frac{1}{2}\operatorname{tr}((XY)^2), \;\;\;\;\;\;(X,Y) \in \mathcal{T}^\ast X^0_N\;.
\end{equation}

By applying the previous map we obtain

\begin{align*}
\mathcal{H} &= \frac{1}{2}\operatorname{tr}((XY)^2)=\frac{1}{2}\operatorname{tr}(UU^{-1}XUU^{-1}YUU^{-1}XUU^{-1}Y) = \\
&= \frac{1}{2}\operatorname{tr}(U^{-1}YUU^{-1}XUU^{-1}YUU^{-1}XU) = \frac{1}{2}\operatorname{tr}(QLQL) = \frac{1}{2}\operatorname{tr}((QL)^2)\;.
\end{align*}

Carrying out the calculations, we recover

\begin{equation}
\mathcal{H} = \frac{1}{2} \sum_{i=1}^N q_i^2 p_i^2 + g^2 \sum_{1 \leq i < j \leq N} \frac{q_i q_j}{(q_i - q_j)^2}\;.
\end{equation}

Here $X$ is unitary, so its eigenvalues lie on the unit circle. Then, with the canonical transformation

\begin{equation}
    \begin{cases}
    q_i = e^{2ia\hat q_i} \\
    \hat p_i = 2ia\, q_ip_i
    \end{cases}
\end{equation}

we have $q_i q_j/(q_i - q_j)^2 = -1/(4\sin^2{a(\hat q_i - \hat q_j)})$, and we get

\begin{equation}
\mathcal{H} = -\frac{1}{4a^2}\Big(\frac{1}{2} \sum_{i=1}^N \hat p_i^2 + g^2 \sum_{1 \leq i < j \leq N} \frac{a^2}{\sin^2{a(\hat q_i - \hat q_j)}}\Big) = -\frac{1}{4a^2}\,\mathcal{H}_{tCM}\;.
\end{equation}

The constant factor only rescales time, so the dynamics is that of the trigonometric CM system.

\subsection{Hyperbolic CM (hCM) system}
The hyperbolic CM system derives from the Hamiltonian on $\mathcal{T}^\ast X^0_N$

\begin{equation}
    \mathcal{H} = \frac{1}{2}\operatorname{tr}((XY)^2), \;\;\;\;\;\;(X,Y) \in \mathcal{T}^\ast X^0_N\;.
\end{equation}

which is the same Hamiltonian as in the trigonometric case, but now $X$ is a positive-definite Hermitian matrix. Operating in the same way as above, with the canonical transformation $q_i = e^{2a\hat q_i}$, $\hat p_i = 2a\, q_ip_i$, we have $q_i q_j/(q_i - q_j)^2 = 1/(4\sinh^2{a(\hat q_i - \hat q_j)})$, and we obtain

\begin{equation}
    \mathcal{H} = \frac{1}{4a^2}\Big(\frac{1}{2} \sum_{i=1}^N \hat p_i^2 + g^2 \sum_{1 \leq i < j \leq N} \frac{a^2}{\sinh^2{a(\hat q_i - \hat q_j)}}\Big) = \frac{1}{4a^2}\,\mathcal{H}_{hCM}\;.
\end{equation}

\section{The Olshanetsky-Perelomov (OP) projection method}
Integrating the equations of motion of these systems directly is hard. However, an alternative method of integration has been developed by M. A. Olshanetsky and A. M. Perelomov in \cite{olshanetskyperelomov}. Their projection method amounts to the integration of the auxiliary Hamiltonian systems, from which the solution to the considered ones can be recovered by projection.

\subsection{Rational CM system}
From free motion on $X^0_N$, described by the Hamiltonian

\begin{equation}
    \mathcal{H} = \frac{1}{2}\operatorname{tr}(Y^2), \;\;\;\;\;\;(X,Y) \in \mathcal{T}^\ast X^0_N\;.
\end{equation}

we obtain the equation of motion $\ddot X = 0$, which integrates directly to 

\begin{equation}
    X = At + B\;\;\;\;,\;\;\;\;A, B \in X^0_N\;.
\end{equation}

Given a similarity transformation

\begin{equation}
    X(t) = U(t) Q(t) U^{-1}(t)
\end{equation}

where $U$ is a unitary $N \times N$ matrix, we differentiate with respect to time, obtaining

\begin{equation}
    \dot X = A = \dot UQU^{-1} + U \dot Q U^{-1} + UQ \dot U^{-1}\;
\end{equation}

Since $\dot U^{-1} = - U^{-1} \dot U U^{-1}$,

\begin{equation}
    \dot X = \dot UQU^{-1} + U \dot Q U^{-1} -UQU^{-1}\dot U U^{-1}\;.
\end{equation}

Then, with some algebraic manipulations, 

\[\dot X = UU^{-1}\dot UQU^{-1} + U \dot Q U^{-1} -UQU^{-1}\dot U U^{-1} = 
    U(U^{-1} \dot U Q + \dot Q - QU^{-1} \dot U )U^{-1}\;.\]

Let $M = U^{-1}\dot U$, then

\begin{equation}
    \dot X = U(MQ + \dot Q - QM)U^{-1} = U(\dot Q - [Q, M])U^{-1} = A\;.
\end{equation}

Now, let $L = \dot Q - [Q,M]$, the previous expression becomes

\begin{equation}
    \dot X = ULU^{-1} = A.
\end{equation}

Differentiating again with respect to time, we get

\begin{equation*}
\begin{split}
\ddot X = 0 &= \dot ULU^{-1} + U\dot L U^{-1} + UL\dot U^{-1} =\\
 &= UU^{-1}\dot ULU^{-1} + U\dot L U^{-1} - ULU^{-1}\dot U U^{-1} =\\
&= U(ML + \dot L- LM)U^{-1}\;.
\end{split}
\end{equation*}

Hence, 

\begin{equation}
    \ddot X = U(\dot L - [L,M])U^{-1} = 0\;.
\end{equation}

We observe that, given an $M$, the matrix $U(t)$ is defined up to initial value by $M = U^{-1}\dot U \iff \dot U = UM$. Setting $U(0) = \mathbb{1}_N$, the previous relations allow us to determine $A$ and $B$:

\begin{equation}
    \dot X = A = U(t)L(t)U^{-1}(t) = U(0)L(0)U^{-1}(0) = \mathbb{1}_N L(0) \mathbb{1}_N = L(0)\;,
\end{equation}

\begin{equation}
    X(0) = B = U(0)Q(0)U^{-1}(0) = \mathbb{1}_N Q(0) \mathbb{1}_N = Q(0)\;.
\end{equation}

Consequently

\begin{equation}
    X(t) = At + B = L(0)t + Q(0)\;,
\end{equation}

completely determining $X$ as

\begin{equation}
    X_{nm} (t) = \delta_{nm} (p_n(0)t + q_n(0)) + (1 - \delta_{nm}) ig \frac{t}{q_n(0) - q_m(0)}\;,
\end{equation}

that represents the solution to the equation of motion $\ddot X = 0$ on $X^0_N$. Finally, the time evolution of $Q(t)$ is given by

\begin{equation}
    Q(t) = U^{-1}(t) X(t) U(t)\;.
\end{equation}

\subsection{Rational CM system with external harmonic potential}
From harmonic motion on $X^0_N$, described by the Hamiltonian

\begin{equation}
    \mathcal{H} = \frac{1}{2}\operatorname{tr}(Y^2 + \omega^2 X^2), \;\;\;\;\;\;(X,Y) \in \mathcal{T}^\ast X^0_N\;.
\end{equation}

we obtain the equation of motion $\ddot X + \omega^2 X= 0$, which integrates directly to 

\begin{equation}
    X = \frac{A}{\omega}\sin{\omega t} + B\cos{\omega t}\;\;\;\;,\;\;\;\;A, B \in X^0_N\;.
\end{equation}

Given the same similarity transformation as above, setting $U(0) = \mathbb{1}_N$ implies $B = X(0) = Q(0)$ and $A = \dot X(0) = L(0)$. Then, we can completely determine $X$ as

\begin{equation}
    X_{nm} (t) = \delta_{nm} (p_n(0)\frac{\sin{\omega t}}{\omega} + q_n(0)\cos{\omega t}) + (1 - \delta_{nm}) ig \frac{\sin{\omega t}}{\omega(q_n(0) - q_m(0))}\;.
\end{equation}

The time evolution of $Q(t)$ is given as already said by

\begin{equation}
    Q(t) = U^{-1}(t) X(t) U(t)\;.
\end{equation}

\subsection{Trigonometric CM system}
Consider now the motion on $X^0_N$ described by the Hamiltonian

\begin{equation}
    \mathcal{H} = \frac{1}{2}\operatorname{tr}((XY)^2)\;, \;\;\;\;\;\;(X,Y) \in \mathcal{T}^\ast X^0_N\;.
\end{equation}

Hamilton's equations read

\begin{equation}
    \dot X = XYX\;\;\;\;,\;\;\;\;\dot Y = -YXY\;,
\end{equation}

so the matrix $C \coloneqq XY$ is conserved:

\begin{equation}
    \dot C = \dot X Y + X \dot Y = XYXY - XYXY = 0\;.
\end{equation}

Hence $\dot X = CX$, which integrates to $X(t) = e^{Ct}X(0)$. Setting $U(0) = \mathbb{1}_N$ as above, we have $X(0) = Q(0)$ and $Y(0) = L(0)$, so

\begin{equation}
    X(t) = \exp{\big(Q(0)L(0)\,t\big)}\,Q(0)\;.
\end{equation}

The eigenvalues of $X(t)$ are $q_i(t) = e^{2ia\hat q_i(t)}$, which give the positions $\hat q_i(t)$ of the trigonometric CM system (up to the time rescaling found above).

\subsection{Hyperbolic CM system}
The hyperbolic case is obtained in the same way, with $X$ positive-definite Hermitian. The solution is again

\begin{equation}
    X(t) = \exp{\big(Q(0)L(0)\,t\big)}\,Q(0)\;,
\end{equation}

and now the eigenvalues of $X(t)$ are $q_i(t) = e^{2a\hat q_i(t)}$.

\section{Lax representations of CM systems}\label{laxcouples}
A Lax pair for the rational CM system was first found by Moser \cite{moser1975}. Calogero then derived Lax representations for CM systems from a general ansatz (see \cite{calogero2003classical}). It turns out that a convenient representation for the systems in analysis is the following:

\begin{subequations}
\begin{align}[left=\empheqlbrace]
L_{nm} &= \delta_{nm} p_n + (1 - \delta_{nm}) \alpha (q_n - q_m) \\
M_{nm} &= \delta_{nm} \sum_{l \neq n}^N \beta(q_n - q_l) + (1-\delta_{nm}) \gamma (q_n - q_m)\;.
\end{align}
\end{subequations}

where $\alpha(q), \beta(q), \gamma(q)$ are functions to be determined. Inserting this ansatz in the Lax evolution equation $\dot L = [L,M]$, and making some simple mathematical assumptions,

\begin{itemize}
    \item $p_n = \dot q_n$,
    \item $\gamma (q) = \alpha'(q)$,
    \item $V(q) = \alpha (q) \alpha (-q)$
\end{itemize}

allow us to derive a dynamical system described by the Lax pair with canonical Poisson brackets $\{q_i, p_j\} = \delta_{ij}$ and Hamiltonian

\begin{equation*}
    \mathcal{H} = \frac{1}{2}\sum_{i = 1}^N p_i^2 + \sum_{1 \leq i < j \leq N} V(q_i - q_j)\;,
\end{equation*}

where $V(q)$ is the interaction potential. Note that here $V$ includes the coupling constant, e.g. $V(q) = g^2/q^2$ in the rational case. Finally, assuming

\begin{itemize}
    \item $\beta(q) = \beta(-q)$,
    \item $\alpha(q_n - q_m)[\beta(q_m - q_l) - \beta(q_n - q_l)] + \alpha(q_n - q_l)\gamma(q_l - q_m) - \alpha(q_l - q_m)\gamma(q_n - q_l) = 0$,
    \item $x = q_n - q_l$ and $y = q_l - q_m$,
\end{itemize}

we get the functional equation

\begin{equation}\label{laxfunctionalequations}
    \frac{\alpha(x) \alpha'(y) - \alpha(y) \alpha'(x)}{\alpha(x+y)} = \beta(x) - \beta(y),  
\end{equation}

Particular solutions to the equation are

\begin{equation}
    \alpha(x) = \frac{ig}{x}\;\;\;\;,\;\;\;\;\beta(x) = \frac{ig}{x^2}\;,
\end{equation}

that give rise to the rational Calogero-Moser system, or 

\begin{equation}
    \alpha(x) = \frac{iga}{\sin{(ax)}}\;\;\;\;,\;\;\;\;\beta(x) = \frac{iga^2}{\sin^2{(ax)}}\;,
\end{equation}

\begin{equation}
    \alpha(x) = \frac{iga}{\sinh{(ax)}}\;\;\;\;,\;\;\;\;\beta(x) = \frac{iga^2}{\sinh^2{(ax)}}\;,
\end{equation}

which correspond respectively to the trigonometric and hyperbolic CM systems. We report for all the three systems considered their Lax representation below.

\subsection{Rational CM system}
The isolated CM system with $V(x) = \frac{1}{x^2}$ can be represented as

\begin{subequations}
\begin{align}[left=\empheqlbrace]
L_{nm} &= \delta_{nm} p_n + (1 - \delta_{nm}) \frac{ig}{q_n - q_m} \\
M_{nm} &= \delta_{nm} \sum_{l \neq n}^N \frac{ig}{(q_n - q_l)^2} - (1-\delta_{nm}) \frac{ig}{(q_n - q_m)^2}\;.
\end{align}
\end{subequations}

\begin{obs}
    Note that for the rational case with external harmonic potential, $L$ and $M$ are not a proper Lax pair for the reduced system, since they satisfy
    
    \begin{equation}
        \dot L - [L,M] = - \omega^2 Q \;.
    \end{equation}
\end{obs}

\subsection{Trigonometric CM system}
Analogously the isolated CM system with $V(x) = {a^2}/{\sin^2{ax}}$ can be represented as

\begin{subequations}
\begin{align}[left=\empheqlbrace]
L_{nm} &= \delta_{nm} p_n + (1 - \delta_{nm}) iga\frac{1}{\sin{a(q_n - q_m)}} \\
M_{nm} &= \delta_{nm} \sum_{l \neq n}^N \frac{iga^2}{\sin^2{a(q_n - q_l)}} \nonumber\\ &\quad - (1-\delta_{nm}) \frac{iga^2}{\sin^2{a(q_n - q_m)}} \cos{a(q_n - q_m)}\;.
\end{align}
\end{subequations}

\subsection{Hyperbolic CM system}
Finally the isolated CM system with $V(x) = {a^2}/{\sinh^2{ax}}$ can be represented as

\begin{subequations}
\begin{align}[left=\empheqlbrace]
L_{nm} &= \delta_{nm} p_n + (1 - \delta_{nm}) iga\frac{1}{\sinh{a(q_n - q_m)}} \\
M_{nm} &= \delta_{nm} \sum_{l \neq n}^N \frac{iga^2}{\sinh^2{a(q_n - q_l)}} \nonumber\\ &\quad - (1-\delta_{nm}) \frac{iga^2}{\sinh^2{a(q_n - q_m)}} \cosh{a(q_n - q_m)}\;.
\end{align}
\end{subequations}

\section{Bäcklund transformations of CM systems}
Finally, let us introduce Bäcklund transformations (BTs) of Calogero-Moser systems. \\
Bäcklund transformations of the CM system were proposed for the first time by Wojciechowski in \cite{wojciechowski1982}. More recently M. Kulkarni and A. Polychronakos in \cite{kulkarni2017} propose a general first-order formulation of the Calogero model in external potentials in terms of a generating function, which simplifies the derivation of its Bäcklund transformations. In this work, we will present BTs in a single form for systems I, II, III and $\textnormal{I}^+$ (see \cite{philip2019soliton} for the proof).

\begin{thm}
Consider $N$ complex canonical coordinates $x_j = q_j, \dot x_j = p_j \in \mathbb{C}$ and $M$ canonical coordinates $z_n, \dot z_n \in \mathbb{C}$. Then, the Bäcklund transformations for CM systems of type $\textnormal{I,II,III}$ and $\textnormal{I}^+$ take the form

\begin{subequations}
\begin{align}[left=\empheqlbrace]
\dot x_j &= - \sum_{k=1,\;k \neq j}^N \alpha(x_j - x_k) + \sum_{n = 1}^M \alpha(x_j - z_n) + \lambda + i\omega x_j\;,\\
\dot z_n &= \sum_{m=1,\;m \neq n}^M \alpha(z_n - z_m) - \sum_{l = 1}^N \alpha(z_n - x_l) + \lambda + i\omega z_n\;,
\end{align}
\end{subequations}

with $\alpha(x)$ the \textit{special Lax function}, with the property $V(x) = - \alpha^2(x)$, $\lambda$ a constant and $\omega = 0$ for all but class $\textnormal{I}^+$.
\end{thm}

The transformations for the cases under analysis are listed below, with $\lambda = 0$.

\subsection{Rational CM system}
Bäcklund transformations of the isolated CM system with $V(x) = {1}/{x^2}$ assume the form

\begin{subequations}
\begin{align}[left=\empheqlbrace]
\dot x_j &= -ig \sum_{k=1,\;k \neq j}^N \frac{1}{x_j - x_k} + ig\sum_{n = 1}^M \frac{1}{x_j - z_n}\;,\\
\dot z_n &= ig \sum_{m=1,\;m \neq n}^M \frac{1}{z_n - z_m} - ig\sum_{l = 1}^N \frac{1}{z_n - x_l}\;.
\end{align}
\end{subequations}

\subsection{Rational CM system with external harmonic potential}
Bäcklund transformations of the  CM system with $V(x) = {1}/{x^2}$ and external potential $W(x) = \frac{1}{2}\omega^2 x^2$ assume the form

\begin{subequations}
\begin{align}[left=\empheqlbrace]
\dot x_j &= -ig \sum_{k=1,\;k \neq j}^N \frac{1}{x_j - x_k} + ig\sum_{n = 1}^M \frac{1}{x_j - z_n} + i\omega x_j\\
\dot z_n &= ig \sum_{m=1,\;m \neq n}^M \frac{1}{z_n - z_m} - ig\sum_{l = 1}^N \frac{1}{z_n - x_l} + i\omega z_n.
\end{align}
\end{subequations}

\subsection{Trigonometric CM system}
Analogously the Bäcklund transformations of the isolated CM system with $V(x) = {a^2}/{\sin^2{(ax)}}$ are given by

\begin{subequations}
\begin{align}[left=\empheqlbrace]
\dot x_j &= -iga \sum_{k=1,\;k \neq j}^N \frac{1}{\sin{a(x_j - x_k)}} + iga\sum_{n = 1}^M \frac{1}{\sin{a(x_j - z_n)}}\;, \\
\dot z_n &= iga \sum_{m=1,\;m \neq n}^M \frac{1}{\sin{a(z_n - z_m)}} - iga\sum_{l = 1}^N \frac{1}{\sin{a(z_n - x_l)}}\;.
\end{align}
\end{subequations}

\subsection{Hyperbolic CM system}
Lastly the Bäcklund transformations of the isolated CM system with $V(x) = {a^2}/{\sinh^2{(ax)}}$ are

\begin{subequations}
\begin{align}[left=\empheqlbrace]
\dot x_j &= -iga \sum_{k=1,\;k \neq j}^N \frac{1}{\sinh{a(x_j - x_k)}} + iga\sum_{n = 1}^M \frac{1}{\sinh{a(x_j - z_n)}}\;, \\
\dot z_n &= iga \sum_{m=1,\;m \neq n}^M \frac{1}{\sinh{a(z_n - z_m)}} - iga\sum_{l = 1}^N \frac{1}{\sinh{a(z_n - x_l)}}\;.
\end{align}
\end{subequations}